% SPDX-License-Identifier: Apache-2.0
\begin{frame}[t]{Aus einer Anfrage einen Graphen erzeugen}
\small CONSTRUCT bildet Tripel aus passenden Variablenbindungen.\par\vspace{3mm}
\begin{activity}[id=sparql-construct,input=code,engine=sparql,seed=data/library.ttl,
reference=answers/sparql-construct.rq,open=manual,limit=5,height=.61\textheight]
\prompt{Erzeugen Sie \texttt{?member :borrows ?book} für jede offene Ausleihe. Nutzen Sie \texttt{:member}, \texttt{:book} und \texttt{:status :open}. Das Ergebnis soll ein Graph sein, keine SELECT-Tabelle.}
\end{activity}
\note{Erwartet: :ada :borrows :b101; :ben :borrows :b102; :cara :borrows :b104. Vertauschen von Subjekt und Objekt verändert den Graphen. Die Ausgangsdaten werden nicht verändert.}
\end{frame}
